\documentclass[runningheads]{llncs}

\usepackage[T1]{fontenc}
\usepackage[utf8]{inputenc}
\usepackage{graphicx}
\usepackage{booktabs}
\usepackage{amsmath}
\usepackage[hyphens]{url}
\usepackage{textcomp}
\usepackage{listings}
\lstset{
  basicstyle=\small\ttfamily,
  breaklines=true,
  frame=single,
  xleftmargin=1em,
  xrightmargin=1em,
}

\begin{document}

\title{Fatawa AI: Evaluating LLM Reliability and Accuracy for Islamic Jurisprudence Queries}
\titlerunning{Fatawa AI: Evaluating LLM Reliability for Islamic Jurisprudence}

\author{Muwaffaq Imam\inst{1} \and
Nursultan Askarbekuly\inst{1} \and
Evgenii Zuev\inst{1}}
\authorrunning{M. Imam et al.}

\institute{Innopolis University, Innopolis, Russia\\
\email{\{m.imam,n.askarbekuly\}@innopolis.university}\\
\email{e.zuev@innopolis.ru}}

\maketitle

\begin{abstract}
Large Language Models (LLMs) are actively used to generate advice and informative content across sensitive domains, including religious guidance. However, their reliability in delivering authoritative, nuanced Islamic Jurisprudence (Fiqh) rulings within a specific school remains a critical vulnerability that demands rigorous, empirical evaluation. In this paper, we present an evaluation of state-of-the-art LLMs---specifically Gemini~2.5 Flash and GPT-5.1---using a newly constructed, expert-validated 298-question benchmark strictly based on the Hanafi school of jurisprudence. By challenging these models with contextual, scenario-based queries, our study reveals substantial gaps in current model reasoning capabilities. We establish a baseline for LLM reliability on Hanafi Fiqh, identifying specific areas where generative AI critically fails.

\keywords{Large Language Models \and Islamic Jurisprudence \and Hanafi Fiqh \and LLM Evaluation \and Hallucination \and Religious AI.}
\end{abstract}

\section{Introduction}

The rapid spread of Artificial Intelligence, particularly Large Language Models (LLMs), presents a blend of opportunities and ethical challenges to Islamic scholarship~\cite{patel2023building,atif2025fiqhqa,islamicEval2025}. The first attempts to explore AI have been already made in the context of Qur'anic retrieval, Hadith question answering, Islamic text search, and religious-domain chatbots~\cite{malhas2020ayatec,malhas2023quranqa,wiharja2022questions,rizqullah2023qasina,mohammed2022english}. Another frontier that lies ahead is the computational generation of religious rulings (fatwa, pl. fatawa).

In Sunni Islam, religious rulings are addressed within the science of jurisprudence (fiqh), that has been traditionally developed by four major schools of law (madhhab, pl. madhahib). The schools extract and document rulings from the two primary sources: the Quran and Prophetic practice (Sunnah), following the methodology established by the early generations of Muslim scholarship.

In light of AI, the broader scholarly community frequently warns against the danger of unguided AI in Fatawa writing~\cite{whitethread2025danger,patel2023building}. An AI system, by its mathematical nature, lacks the accountability and moral agency required to fulfil the role of a traditional Mufti. Consequently, to quantify these deficiencies, new benchmarks have progressively emerged, such as IslamicLegalBench~\cite{elmahjub2026islamiclegalbench} and FiqhQA~\cite{atif2025fiqhqa}. These benchmarks attempt to evaluate LLMs on a spectrum of Islamic jurisprudence knowledge, moving beyond simple factual recall into the realm of legal reasoning.

Building upon these foundational efforts, our study shifts focus to a targeted evaluation: we investigate whether modern, commercially available LLMs, when acting strictly on their pre-trained weights without external retrieval, can be trusted to provide reliable and accurate answers to questions in the domain of Islamic jurisprudence. We specifically target the Hanafi madhhab, one of the four major schools of law in Sunni Islam. To conduct the study and facilitate further research we also release a 298-question benchmark across five categories, with reference answers in accordance with the Hanafi school of law.

\section{Context and Prior Work}

\subsection{LLMs in Fatawa Generation}
The idea of utilising Large Language Models to autonomously generate fatwas (Islamic legal rulings) is a controversial topic. Prominent Muftis and Islamic scholars warn against the intrinsic dangers of utilising AI as a standalone, primary source for religious rulings~\cite{whitethread2025danger}. The core of this concern lies in the realisation that AI systems inherently lack moral accountability, operate in a vacuum devoid of contextual awareness, and lack the grounded reasoning capacity necessary to interpret rulings in alignment with the higher objectives of Islamic law (Maqasid al-Shari'ah).

Furthermore, LLMs exhibit a dangerous propensity to ``hallucinate''~\cite{lin2021truthfulqa,islamicEval2025}, partially owing to their very nature as probabilistic token generators and also issues in the pre-training data. In a religious context, a hallucination is not merely an error; it can manifest as a fabricated Arabic citation attributed to a canonical scholar, or the dangerous conflation of distinct, sometimes opposing, schools of jurisprudence.

Giving rulings can be roughly categorized into answering known questions and answering newly arisen problems that have not been previously addressed. Answering known questions is fundamentally a problem of knowledge retrieval, and these should be possible to solve if the right pre-training data is provided and hallucination issues are addressed. While answering questions about newly arisen problems is a significantly more challenging task that requires applying the rigorous methodology of deriving rulings from the first principles of religion, the process known as \textit{Ijtihad}, and it is not clear whether LLM models will ever be able to reach that level of sophistication given an ongoing debate whether LLMs can actually apply reasoning and produce knowledge \cite{illusion-of-thinking}.

\subsection{Existing Benchmarks}
Earlier Islamic NLP resources have focused mainly on retrieval and question answering over religious texts. Qur'anic QA benchmarks include annotated Al-Qur'an QA corpora, verse-level retrieval collections, and shared tasks over passage retrieval and reading comprehension~\cite{alqahtani2018annotated,malhas2020ayatec,malhas2023quranqa}. Related work has also explored Hadith and Sirah-based QA~\cite{wiharja2022questions,rizqullah2023qasina}, non-factoid Islamic text QA~\cite{qamar2024benchmark}, and Islamic-domain chatbot datasets~\cite{mohammed2022english}. These resources establish the value of religious-domain NLP, while broader Arabic benchmarks such as ArabicMMLU evaluate general multitask competence rather than Fiqh-specific ruling reliability~\cite{koto2024arabicmmlu}.

To address these shortcomings systematically, the academic NLP community has established targeted initiatives. The IslamicEval 2025 shared task was explicitly aimed at detecting and mitigating hallucinations in Islamic text generation~\cite{islamicEval2025}. Multi-school benchmarks such as IslamicLegalBench~\cite{elmahjub2026islamiclegalbench} and FiqhQA~\cite{atif2025fiqhqa} have emerged to quantify these failures. IslamicLegalBench revealed that state-of-the-art LLMs exhibit significant limitations, hallucinating frequently and failing to achieve high accuracy on complex legal reasoning tasks.

FiqhQA~\cite{atif2025fiqhqa} evaluated LLMs across all four Sunni schools over 120 questions and introduces the concept of abstention evaluation, a reliability dimension also studied in broader work on when LLMs should decline to answer~\cite{madhusudhan2024llms}. In FiqhQA, GPT-4o achieved its highest English zero-shot correctness on Hanafi questions, a result the authors speculate may reflect the greater availability of Hanafi material in public web data; however, this pattern was not uniform across all models or languages.

Unlike any of the previous papers, our work focuses exclusively on the Hanafi school. We also contribute an expert-validated, category-balanced benchmark of 298 Hanafi questions to facilitate further research into knowledge retrieval, abstention ability, and accuracy of Islamic question-answering systems.

\section{Fatawa AI: Dataset and Methodology}

\subsection{Dataset Construction and Validation}
Our benchmark is a highly structured compilation of 298 Question--Answer (Q\&A) pairs, exclusively representing rulings from the Hanafi school. It was derived from a larger cleaned source pool of approximately 10{,}000 fatawa retrieved from a legacy fatwa database through its API. The raw JSON returned by the database was cleaned, normalised into tabular records, and prepared for classification and expert review.

The cleaned source pool was then organised into a data taxonomy of 10 top-level categories and 92 highly specific sub-categories. Initial category titles were produced with an LLM- and embedding-assisted classification pipeline: each candidate question and its corresponding answer were embedded and matched against vector representations of the candidate sub-categories using cosine similarity. These suggested classifications were then validated by the domain expert and review team, who corrected category assignments where necessary. The validated taxonomy was used to organise the candidate pool and select a balanced benchmark covering daily life scenarios---ranging from modern financial transactions to classical purification rituals. Within this paper, we release the first 298 question-answer pairs from the dataset.

\subsection{Evaluation Methodology}

\subsubsection{Test Set Construction.}
For the evaluation, we first selected the five most frequent categories that each appeared at least 500 times in the cleaned source pool. We then sampled 60 questions per selected category, yielding a provisional test set of 300 questions:
\textit{Ibadat} (acts of worship), \textit{Zawaj} (marriage, divorce, expenditures), \textit{Mu'amalat} (transactions and trusts), \textit{Hazr} (prohibition, dress, adornment), and \textit{Dhikr} (remembrance, purification, ethics).
Two questions were excluded from the final paired analysis because the experiment export contained incomplete or duplicated model-output pairings from a manual recovery step. The final evaluation therefore contains 298 paired questions and 596 model-answer judgements.

\subsubsection{Queried Models.}
Both models were queried in zero-shot mode---no retrieval context was provided, so results reflect purely what each model has internalised during pre-training. The same system prompt was applied to both models:

\begin{lstlisting}
You are an expert Islamic scholar specializing in Fiqh,
specifically adhering to the Hanafi Madhab.
Answer user inquiries in Arabic (al-Fusha or accessible
standard Arabic).
Strictly cite rulings based on the Hanafi school. Do not
offer opinions from other schools (Maliki, Shafi'i, or
Hanbali) unless explicitly asked for a comparison.
Tone: Respectful, authoritative, and concise.
Task: Analyze the user's question and provide a direct
Hanafi ruling.
\end{lstlisting}

The user turn was the template \texttt{"user~question~is:~\{query\}"}. Both models used their reasoning capabilities: Gemini~2.5 Flash with a dynamic thinking budget and GPT-5.1 with low reasoning effort.

\subsubsection{LLM-as-a-Judge.}
Given that a correct Fiqh ruling can be phrased in many semantically equivalent ways, string-matching metrics such as BLEU or Exact Match are unsuitable. Prior work on semantic and model-based evaluation similarly shows the need for meaning-sensitive judgement beyond surface overlap~\cite{zhang2019bertscore,yuan2021bartscore,zhong2022towards,liu2023geval,zheng2023judging}. We therefore used an LLM-as-a-judge protocol, comparing each generated answer against the domain expert's ground-truth answer. After human validation (Section~\ref{sec:judge-validation}), the final automated judge was Grok~4.3 via OpenRouter with \texttt{reasoning\_effort=medium}, \texttt{temperature=0}, and strict JSON-schema output.

The judge prompt defined the reference answer as authoritative and instructed the model not to overrule it using outside Fiqh knowledge. The judge was also instructed to ignore differences in wording, length, style, and citation format unless they changed the legal ruling. It returned a single JSON object containing one score, \texttt{\{"rating": -1/0/+1\}}, on the following ordinal scale:

\begin{itemize}
  \item \textbf{$+1$ (Identical Meaning).} The generated answer is semantically equivalent to the ground truth in intent, condition, and ruling. A safely deployable response.
  \item \textbf{$0$ (Partially Correct).} The model agrees in some fundamental aspects but hallucinates conditions, misses critical exceptions, or blends rulings from another school.
  \item \textbf{$-1$ (Contradictory).} The model provides an answer that is fundamentally opposite to the ground truth (e.g., stating an act is \textit{Haram} when the Hanafi ruling is \textit{Halal}). The most dangerous failure mode.
\end{itemize}

\subsubsection{Judge Validation.}
\label{sec:judge-validation}
To assess the reliability of the LLM judge, we conducted a human validation study on a stratified sample of 30 unique questions (6 per category), evaluated for both model outputs, yielding 60 judge--human comparison pairs. The domain expert independently assigned a $\{-1, 0, +1\}$ score to each (question, ground-truth, model-answer) triple without seeing the automated judge score.

Because the score labels form an ordered scale, we used ordinal Krippendorff's $\alpha$ as the primary reliability statistic~\cite{krippendorff2011computing}. This metric penalises severe disagreements ($-1$ vs $+1$) more than adjacent disagreements ($0$ vs $+1$), and unlike Cohen's $\kappa$, it naturally extends to future validation with more than two raters or missing ratings. We also report linear-weighted Cohen's $\kappa$ as a secondary two-rater robustness check~\cite{cohen1968weighted}.

We did not use GPT-5.1 or Gemini~2.5 Flash as judges because they were the systems under evaluation. Instead, we first tested a Claude judge, but its agreement with the human expert was weak on the ordinal reliability metric. We then evaluated Grok~4.3 under low and medium reasoning effort. Table~\ref{tab:judge-validation} reports the resulting judge--human agreement.

\begin{table}[t]
\centering
\caption{Human validation of candidate LLM judges over 60 model-answer pairs.}
\label{tab:judge-validation}
\setlength{\tabcolsep}{4pt}
\footnotesize
\begin{tabular}{lccc}
\toprule
\textbf{Candidate judge} & \textbf{Exact agree.} & \textbf{Linear $\kappa$} & \textbf{Ordinal $\alpha$} \\
\midrule
Claude judge       & 66.7\% & 0.432 & 0.491 \\
Grok~4.3 low       & 76.7\% & 0.570 & 0.651 \\
Grok~4.3 medium    & 71.7\% & 0.535 & \textbf{0.702} \\
\bottomrule
\end{tabular}
\end{table}

Although Grok~4.3 with low reasoning achieved the highest exact agreement and weighted $\kappa$, Grok~4.3 with medium reasoning achieved the strongest ordinal Krippendorff's $\alpha$ ($0.702$), our primary metric. We therefore used Grok~4.3 with \texttt{reasoning\_effort=medium} for the final automated evaluation.

\section{Results}

We evaluate models primarily through \emph{Simple Accuracy} (percentage of $+1$ judgements), as a partially correct Fiqh ruling often lacks the precision to be safely acted upon by a layman.

\begin{equation}
\text{Accuracy} = \frac{\text{Count of }+1\text{ judgements}}{N} \times 100\%
\end{equation}

\subsection{Overall Performance}
Table~\ref{tab:overall} presents aggregate performance across $N=298$ questions per model. GPT-5.1 achieved an overall accuracy of 53.36\%, slightly outperforming Gemini~2.5-Flash (52.35\%) by 1.01 percentage points. This difference was not statistically significant under McNemar's paired test ($p = 0.8174$, $\chi^2 = 0.0533$).

Both models fall drastically short of reliability thresholds accepted in legal or medical AI applications ($>95\%$). The model comparison should be interpreted with caution: the one-point gap is not meaningful evidence of a consistent architectural advantage, and both systems remain unsafe for direct legal-religious deployment.

\begin{table}[t]
\centering
\caption{Overall Simple Accuracy and score distribution across 298 paired Hanafi Fiqh instances per model.}
\label{tab:overall}
\setlength{\tabcolsep}{4pt}
\footnotesize
\begin{tabular}{lccccc}
\toprule
\textbf{Model} & \textbf{N} & \textbf{$+1$ Correct} & \textbf{$0$ Partial} & \textbf{$-1$ Contradict.} & \textbf{Acc.} \\
\midrule
GPT-5.1          & 298 & 159 (53.4\%) & 83 (27.9\%) & 56 (18.8\%) & 53.36\% \\
Gemini-2.5-Flash & 298 & 156 (52.3\%) & 71 (23.8\%) & 71 (23.8\%) & 52.35\% \\
\bottomrule
\end{tabular}
\end{table}

\subsection{Score Distribution and Failure Mode Analysis}
Table~\ref{tab:overall} separates errors into partially correct ($0$) and contradictory ($-1$) responses. This distinction matters for deployment: partial errors indicate incomplete knowledge, while contradictory errors ($-1$) represent active misinformation. A model with a high $-1$ rate should be treated as actively unsafe for direct deployment, regardless of overall accuracy.

\subsection{Cross-Category Performance}
Table~\ref{tab:categories} shows per-category accuracy. The nature of the context overwhelmingly dictates model reliability, confirming that Fiqh is not a monolithic database of facts but a diverse landscape of reasoning skills.

Descriptively, Gemini~2.5-Flash achieved higher accuracy in \textit{Dhikr} (ethics and purification), \textit{Ibadat} (acts of worship), and \textit{Mu'amalat} (transactions and trusts), while GPT-5.1 achieved higher accuracy in \textit{Zawaj} (family law, marriage, divorce) and \textit{Hazr} (prohibition, dress, adornment). These category-level reversals indicate that aggregate accuracy obscures distinct failure profiles across branches of Fiqh.

To assess whether these apparent differences were reliable, we applied exact McNemar tests within each category and Holm correction across the five category-level comparisons. Only the \textit{Hazr} gap remained significant after correction (GPT-5.1 60.00\% vs. Gemini~2.5-Flash 40.00\%; raw $p=0.0018$, Holm-adjusted $p=0.0092$). The \textit{Ibadat} gap favored Gemini~2.5-Flash before correction (51.67\% vs. 35.00\%; raw $p=0.0414$) but did not survive Holm correction (adjusted $p=0.1656$); all other category-level model differences were non-significant. Chi-square tests of model accuracy by category further showed significant variation for GPT-5.1 ($\chi^2=13.63$, $p=0.0086$), but not for Gemini~2.5-Flash ($\chi^2=5.09$, $p=0.278$).

\begin{table}[t]
\centering
\caption{Cross-category accuracy comparison.}
\label{tab:categories}
\setlength{\tabcolsep}{3pt}
\footnotesize
\begin{tabular}{p{1.7cm} p{3.0cm} p{2.1cm} p{2.4cm} p{1.4cm}}
\toprule
\textbf{Category} & \textbf{English} & \textbf{GPT-5.1} & \textbf{Gemini-2.5-Flash} & \textbf{Winner} \\
\midrule
Mu'amalat & Transactions and Trusts           & 50.85\% (30/59)          & 55.93\% (33/59)          & Gemini \\
Hazr      & Prohibition, Dress, Adornment     & 60.00\% (36/60)          & 40.00\% (24/60)          & GPT-5.1 \\
Dhikr     & Remembrance, Purification, Ethics & 54.24\% (32/59)          & 57.63\% (34/59)          & Gemini \\
Zawaj     & Marriage, Divorce, Expenditures   & \textbf{66.67\%} (40/60) & 56.67\% (34/60)          & GPT-5.1 \\
Ibadat    & Acts of Worship                   & 35.00\% (21/60)          & 51.67\% (31/60)          & Gemini \\
\bottomrule
\end{tabular}
\end{table}

\section{Discussion and Ethical Implications}

The data generated from this evaluation carries profound implications for the deployment of AI in religious contexts. An overall accuracy hovering only in the low-50\% range confirms that querying a base LLM for Islamic legal advice is currently an unsafe practice. The fundamental flaw observed is not just the lack of knowledge, but the confident presentation of incorrect knowledge.

This failure makes abstention ability central to any safe religious-domain system, confirming the motivation behind FiqhQA's abstention-oriented evaluation~\cite{atif2025fiqhqa} and broader work on teaching LLMs when to decline answering~\cite{madhusudhan2024llms}. A suitable scoring methodology for this domain should reward only fully correct answers, assign a neutral score to explicit abstention, and penalise incorrect answers with a negative score. This logic is relevant not only for post-hoc evaluation by researchers, but also for training and alignment: reward functions that treat plausible but wrong religious rulings as merely low-quality answers may actively reinforce unsafe behaviour.

Under such a reward-penalty schema, partially correct answers should generally be treated as incorrect for deployment purposes. A partially correct Fiqh answer may still omit a decisive condition, introduce a false exception, or combine rulings from different schools; in a sensitive domain, users need correctness with appropriate certainty rather than approximate semantic overlap. This creates a deeper challenge for LLMs: it is not clear that they can reliably estimate the probability that their own generated ruling is true, since they lack self-awareness and do not have direct introspective access to the truth conditions of their outputs. Reliable abstention may therefore require external grounding, calibrated uncertainty estimation, and explicit refusal mechanisms rather than confidence expressed by the model's surface fluency.

When an LLM scores $0$ (partially correct) or $-1$ (contradictory), it rarely signals uncertainty to the user. Instead, it generates an authoritative-sounding Fatwa, often formatted in classical Arabic, which mimics the structure of an authentic scholarly ruling. For an average layman, distinguishing between a deeply grounded Hanafi ruling and a beautifully articulated hallucination is practically impossible.

One mitigation is to require generated answers to be grounded in verifiable references to authoritative jurisprudence source material rather than produced solely from parametric memory. Here, Retrieval-Augmented Generation and appropriate training data can play a central role by constraining the model to retrieved passages from vetted Hanafi texts. However, citations alone are not sufficient: systems must also provide transparent mechanisms for verifying that quoted passages, page references, and links actually correspond to the cited source, since fabricated citations or fake links can create a false sense of scholarly reliability~\cite{niu2023ragtruth,song2024rag}.

Overall, the results directly support the urgent warnings from the traditional scholarly community. The deployment of ungrounded LLMs in this space risks widespread dissemination of religious inaccuracies, which can severely impact individual observances and family law implementations. Therefore, the ``zero-shot'' or ``direct-query'' method for religious rulings should be actively discouraged by developers building Islamic LLM applications.

\section{Conclusion and Future Work}

Our findings empirically demonstrate that while Large Language Models possess the raw potential to understand and process jurisprudence questions, they currently fall far short of expert-level reliability. Even in the best-performing category, GPT-5.1 reached only 66.67\% accuracy on family-law questions; at the aggregate level, both models remained close to 52--53\%. Base models therefore cannot be trusted as standalone Muftis. Their internal weights represent an amalgamation of conflicting data rather than a structured knowledge hierarchy of a single Fiqh school.

Future work should expand the human validation beyond a single expert rater, and evaluate fine-tuned Arabic LLMs such as ALLaM and Fanar to assess whether domain-specific Arabic pre-training narrows these gaps~\cite{bari2024allam,team2025fanar}. Most importantly, our future roadmap includes comparing these baseline parametric models against Retrieval-Augmented Generation (RAG) systems. Given the catastrophic failures in zero-shot reasoning, we hypothesise that grounding the LLM on pre-authenticated Hanafi texts is the only viable path to safety, while recent work on RAG hallucination shows that retrieval grounding itself must still be evaluated carefully~\cite{niu2023ragtruth,song2024rag}. Early adoption tests where such systems answered over 4{,}000 questions as a safe research assistant to human domain experts lend support to this hypothesis: the future of Fatawa AI lies not in autonomous ruling generation, but in intelligent, validated retrieval at scale.

\section{Data Availability}

The 298-question Hanafi Fiqh benchmark will be released on Hugging Face. The public dataset will contain the question, reference answer, title, category, and subcategory fields only; model outputs, judge prompts, judge scores, and intermediate evaluation artifacts are not included in the public release. The final manuscript version will include the dataset repository URL.

\begin{credits}
\subsubsection{\ackname} The authors thank the human domain expert who curated and validated the 298-question Hanafi Fiqh benchmark.

\subsubsection{\discintname} The authors have no competing interests to declare that are relevant to the content of this article.
\end{credits}

\bibliographystyle{splncs04}
\bibliography{main}

\end{document}
